\section{Applicable Intervention Evidence and Costs}\label{sec:defenses}
\subsection{Which results answer a deployment question?}
The original 43-core snapshot has a complete scoped intervention disposition. Sixteen contain evaluated mitigation components; two use fixed guardrails, four evaluate interventions for other objectives, 13 have no evaluated intervention within the inspected scope, four are diagnostic components, and four have unbound defense assertions. This frozen inventory retains 33 result records across 22 families and 66 utility/cost field positions (Supplement S5). New sources have their own extraction rows (S10); they are not silently pooled into these audit denominators.

For deployment comparisons, we select records in which the source deliberately suppresses the studied channel or detectability and measures either that inference outcome or a cost of the mechanism. Fixed common guardrails, ordinary VPN robustness, access blocking, contention noise, signal-generating privacy filters and channel-only calibration remain contextual records. The companion selector reports every inclusion and retained contextual row. The selected records still mix executed inference, offline/simulated evidence and component costs; they have no common denominator for a defense-failure rate.

Table~\ref{tab:defense-types} organizes eight defense types by the receiving view and the result available for a deployment decision. Combined with the three cases, it identifies what can be used immediately and which additional binding the intended claim requires. Component measurements retain their engineering value even when they do not price an entire service.
\begin{table*}[t]\centering\footnotesize
\caption{Defense types and the evidence available for choosing them. These representative conditions span confidentiality, attribution and human--agent detection; their outcomes are not a common efficacy ranking.}\label{tab:defense-types}
\begin{tabularx}{\textwidth}{@{}p{.16\textwidth}p{.30\textwidth}Y@{}}\toprule
Type & Evaluated view, target and source & Established result and applicable boundary\\\midrule
Padding / batching & Network packet lengths and groups; topic inference. NetEcho and Whisper Leak \cite{netecho2025,whisperleak2025}. & Collected padding traces retain inference signal; packet-group batching is evaluated by replay. A protected upstream length view does not cover downstream re-emission or residual grouping. Live incremental cost remains \NRR{}.\\
Constant-rate / transport cover & Network workflow fingerprint; Dangol's OHTTP testbed \cite{dangol2026workflow}. & The co-located loopback setting yields 0.30 recovery at about $13\times$ bytes and $11\times$ latency. Transformed-capture results are separate; neither result supplies objective agent-task utility.\\
Concurrent decoy agents & Destination/domain observer; private prompt and trait inference. WRA \cite{jeong2026wra}. & Real/decoy requests execute in parallel; interleaved-domain evaluation reduces one-decoy $E_{\rm func}$ from 0.770 to 0.610. Terminal decoy utility/cost are \NRR{}; report usefulness is a separate visibility condition.\\
Action humanization & Mobile touch traces; human--agent detection. AHB \cite{zhu2026turing}. & Executed combined humanization changes Trip Planning success from 0.75 to 0.15. Fake actions can alter state; offline trajectory edits and other online configurations are separate conditions, with different utility implications.\\
Cache isolation / keying & Co-tenant prefix timing; SafeKV and KVGov \cite{chu2025safekv,addagada2026kvgov}. & SafeKV reports post-trigger AUC near 0.5 and separate TTFT/throughput comparisons. KVGov's attack results are simulated; its measured HMAC median $1.3\,\mu$s prices key computation, excluding lost reuse. Each cost retains its mechanism and workload.\\
Cache timing equalization & Co-tenant cache-fusion timing; SpliceDefense \cite{spliceleak2026}. & Constant-Time Boundary Fusion (CTBF) suppresses the tested signal over 15 tasks/75 steps. Quantized Chunk Padding (QCP) adds mean 7.5 tokens and $<5$ms/$<1\%$ prefill time; this does not price combined deployment.\\
Timestamp perturbation & Recorded website action metadata; identity attribution. UI Traces \cite{uitraces2026}. & Fitting on delayed traces recovers substantial predictive performance under offline timestamp edits that retain action types/order. Live execution delay and task outcomes are not measured by the trace transformation.\\
Oblivious / unshared memory & Host or local resource observations; TDXRay and Gao et al.\ \cite{hornetz2026tdxray,gao2025iknow}. & TDXRay measures oblivious-map lookup cost; Gao measures loading/memory cost after disabling zero-copy. These are component measurements, not post-mitigation prompt-recovery or complete serving-cost evaluations.\\\bottomrule
\end{tabularx}
\par\smallskip\raggedright\footnotesize
The table summarizes evaluated conditions or measured components, rather than treating proposed mechanisms as deployed defenses. TDXRay's TEE setting specifies the host adversary; it does not establish that TEE residency alone prevents the studied channel. The result-level inventory retains the separate conditions and missing fields.
\end{table*}


\subsection{What changes when records, rather than papers, are compared?}
We apply the comparison rules to the 33 records, retaining all 528 unordered pairs in the companion audit. The most interpretable subset is the 13 pairs from within the same study family: seven (53.8\%) support a source-qualified descriptive comparison, while six (46.2\%) concern different primitives, outcomes or cost objects. Being reported in one paper therefore does not make two results one operating point. Of the seven aligned-outcome pairs, four use positively different collection/execution populations and three lack a verified matched-test or fitting binding. Their values cannot be subtracted as paired effects. This leaves intact the valid clean/adapted contrasts \emph{within} the UI Traces and Dangol capture records.

Two examples show the practical change. Dangol's capture and live results describe the same workflow-recovery problem but are different operating points. SpliceDefense's CTBF signal result and QCP cost are useful component measurements, not a measured combined frontier. The audit turns these distinctions into explicit decisions with source locators, rather than a warning based only on terminology. Supplement S8 gives the wider, mechanism-conditioned pair inventory and directional transfer judgments; the counts describe our coded comparisons, not mistakes by the original authors.

Table~\ref{tab:six-check-audit} applies the six reporting checks of Section~\ref{sec:agenda} to the 16 evaluated-mitigation families. Targets, receiving views and placement are usually source-bound; deployment-level joins need more care. AHB supplies detectability and objective task success without a measured execution-cost endpoint. Dangol's live result supplies inference and protocol costs without objective task utility. Speculation couples its inference score to token payload, not a wire-cost or task-success measurement. These complementary results identify the additional measurement needed for the intended deployment, rather than imply that each component study failed its own objective.
\begin{table}[t]\centering\footnotesize
\caption{Six checks over 16 evaluated-mitigation families. Counts apply to the source-scoped claim; no composite quality score is computed.}\label{tab:six-check-audit}
\begin{tabularx}{\columnwidth}{@{}Yrrrr@{}}\toprule
Check & Bound & Partial & \NRV{} & \NA{}\\\midrule
1. Target/protection relation &16&0&0&0\\
2. Receiver fields/permissions &13&3&0&0\\
3. Placement/action dispatch &13&2&1&0\\
4. Protected fit/test/access &0&4&5&7\\
5. Tested workloads/shifts &14&2&0&0\\
6. Joint inference/utility/cost &0&16&0&0\\\bottomrule
\end{tabularx}
\par\smallskip\raggedright\footnotesize
Bound: the required item is source-bound for the scoped claim. Partial: supported elements coexist with an unbound element. Check4 includes the observer's training-data access and matching, not just existence of a fitted model. Check6 asks for an end-to-end deployment join, not a requirement that every component study perform one. No whole-check cell is solely \NRR{}; individual unreported subfields remain \NRR{} in the 96-cell inventory.
\end{table}


\subsection{Adaptation, training access and publication runs}
Protected-data fitting (D1) and a paired comparison on the same defended test condition (D2) are separate. In the five multifeature-trace subset, UI Traces, NetEcho and Dangol establish D1; UI Traces and Dangol establish D2. Whisper Leak/AHB exact fit bindings remain unresolved, giving source-status bounds of three--five and two--five respectively. These bounds describe this subset, not all learned attacks. Early Bird's $K$-specific hit/miss calibration is a separate D1 stage before oracle simulation \cite{earlybird2025}.

Labeled protected data must also be obtainable under the observer's permissions. Paired samples, tuning and compute budgets need their own bindings; source-local refitting alone proves neither access nor fairness. For AHB, released code confirms condition-specific refits (D1-code), but the historical Table~1 run (D1-run) and a paired common-test adaptation comparison (D2-paired) remain \NRV{} \cite{zhu2026turing}. The state-change case uses the separately reported task-success result, not an adaptive-detector effect.

\subsection{Costs and excluded quantitative assertions}
Interpret cost as the object measured. Dangol's live 0.30 and transformed-capture 0.17 cannot be subtracted as a paired deployment effect. Its 0.167 timing-feature ablation is not an executed timing masker. Early Bird's 6ms/4\% semantic prototype does not price KV-granularity protection; KVGov's HMAC microcost omits lost cache reuse. SpliceDefense combines Constant-Time Boundary Fusion (CTBF) and Quantized Chunk Padding (QCP); mean 7.5 padding tokens and $<5$ms/$<1$\% describe QCP alone, not the combined cost \cite{dangol2026workflow,earlybird2025,addagada2026kvgov,spliceleak2026}.

Likewise, Speculation measures 8.7$\times$ token payload; Gao measures loading cost and TDXRay map-lookup cost \cite{wei2024speculation,gao2025iknow,hornetz2026tdxray}. SafeKV reports both asynchronous detector latency and serving TTFT/throughput measurements; its five-second post-trigger leakage window and serving comparisons retain their own conditions \cite{chu2025safekv}. These are useful measurements of the named objects, rather than interchangeable prices. Supplements S5/S8 preserve the mode-specific bindings and source corrections.

Four assertions remain outside quantitative comparison because their result/run binding is unresolved: AgentPrint's supplementary analysis, Weiss's vendor confirmation, Gu's provider changes and PromptEcho's 3.2$\times$/1.4$\times$ padding overhead \cite{agentprint2025,weiss2024prompt,gu2025promptcaching,prompt_echo2026}. The three deployment cases use the other, source-bound records.

\paragraph{Positive counterconditions from the update}
CachePrune reports marked-token attack suppression together with question-answering/summarization utility and serving performance under selective KV sharing; contextual reconstruction and detector false negatives remain separate endpoints \cite[Section 5]{wu2026cacheprune}. OTRO's epoch padding reduces access-count information conditional on prompt byte length from 3.94 to 0.04 bits \cite[Section VII]{lee2026otro}. Its main TTFT figures omit padding, reporting about 4.6\% mean overhead for Llama-3.1; a derived upper padding cost of about 2.8ms per epoch is reported separately. The usable positive results are thus selective-cache protection with quality/performance, and conditional tokenizer leakage reduction with separately bounded costs. They preclude reading the old inventory as evidence that no useful joint defense evaluation exists. Reuse concerns the tested mechanism, observer and workload, not every subsequent tool or frontend.
